\documentclass[11pt]{article}
\usepackage{kristiania-workbook}

\begin{document}
\begin{center}
  {\LARGE\bfseries Lock down your VM with a firewall}\\[3pt]
  {\large Session 4 -- Networking \& firewall controls}\\[5pt]
  {\small Exercise 4 \textperiodcentered\ Session 4 lab \textperiodcentered\ Estimated time: 75--100 min}
\end{center}
\vspace{2pt}\hrule\vspace{1.1em}

Write your answers in the answer document \cmd{SKY2100\_Exercise4\_Answers.docx}. Last week you saw what \emph{too much} traffic does;
today you control \emph{which} traffic is allowed at all. You will turn your Ubuntu VM into a
\textbf{default-deny} host with UFW, open only the ports you can justify, and prove the rules work by
testing from the host. This is the same reasoning you will later apply to an Azure Network Security
Group.

\noindent\textbf{Caution:} Keep the SSH port (22) open throughout. If you are working on a remote VM,
closing it can lock you out. Test changes locally first.

\wbsection{1\quad Baseline: what is open now?}
\task{List the ports your VM is currently listening on:} \cmd{ss -tulpn}. Note which services are
bound to all interfaces (\cmd{0.0.0.0} / \cmd{[::]}) versus localhost only.

\task{For each listening port, decide: does it need to be reachable from outside?}

\wbsection{2\quad Configure a default-deny firewall}
\task{Set the policy and open only what you need:}
\begin{quote}\ttfamily
sudo ufw default deny incoming\\
sudo ufw default allow outgoing\\
sudo ufw allow 22/tcp\\
sudo ufw allow 80/tcp\\
sudo ufw enable
\end{quote}

\task{Inspect the result:} \cmd{sudo ufw status verbose}. Which ports did you allow, and why exactly
those?

\wbsection{3\quad Test that the rules actually work}
\task{Start a web server if one is not running} (\cmd{sudo apt install -y nginx}) and find the VM's
address with \cmd{ip a}.

\task{Test from your \emph{host} machine, not from inside the VM} --- UFW does not filter loopback
traffic, so \cmd{localhost} would never show a block. On the host, test an \emph{allowed} and a
\emph{denied} port:
\begin{quote}\ttfamily
curl -I http://<vm-ip>:80\quad(\# allowed -> responds)\\
curl --max-time 5 http://<vm-ip>:8080\quad(\# denied -> times out)
\end{quote}
\note{VirtualBox users: VirtualBox NAT is not reachable from the host --- switch the VM's network
adapter to \emph{Bridged} or \emph{Host-only} first. VMware NAT works as it is.}

\task{For comparison, disable the firewall for a moment} (\cmd{sudo ufw disable}), repeat the second
\cmd{curl}, then \cmd{sudo ufw enable} again. Without the firewall, the closed port answers
``connection refused'' at once.

\task{Briefly, what is the difference between how a \emph{denied} port and a \emph{closed} port fail?
(timeout vs.\ connection refused)}

\task{Re-run \cmd{ss -tulpn} and compare with your baseline.}

\aha{Compare the two \cmd{ss} outputs: they are almost identical. The firewall did not stop a single
service from running; every daemon still listens exactly as before. What changed is who can reach it
from outside. A firewall controls reachability, not whether software runs, and that distinction is the
whole point of the attack-surface idea from the lecture: the service is still there, you have just shut
the door in front of it. The two failures you saw also differ for a reason worth remembering. A closed
port with no firewall in the way answers with a quick ``connection refused'', while a port the firewall
drops simply times out, because the packet is silently discarded and the attacker is left waiting with
no reply.}

\wbsection{4\quad Cloud transfer: Microsoft Learn}
\task{Revisit ``Describe Azure networking services''} and focus on \textbf{Virtual Network (VNet)} and
\textbf{Network Security Group (NSG)}.
\par\noindent\mbox{\href{https://learn.microsoft.com/training/modules/describe-azure-networking-services/}{learn.microsoft.com/training/modules/describe-azure-networking-services}}

\task{Map your UFW rules to an NSG.} For each rule --- \cmd{allow 22/tcp} (SSH), \cmd{allow 80/tcp}
(HTTP) and the default deny --- write in the answer document what the equivalent NSG rule would look
like.

\wbsection{5\quad Azure reflection}
\task{An NSG protects a whole subnet; UFW protects one host. Why is running \emph{both} an example of
defence in depth, rather than redundant effort?}

\wbsection{6\quad Knowledge check}
\begin{enumerate}
\item What does a socket consist of?\\
      \cbox an IP address only \quad \cbox a port only \quad \cbox an IP address + a port \quad
      \cbox a MAC address + a port
\hint{Two things together: which machine, and which program on it.}
\item Why is \textbf{default-deny} the only safe default policy for a firewall?
\hint{You cannot foresee every threat, so you open only what you can justify.}
\item Which port is which? \cmd{22} = \rule{2.2cm}{0.3pt}; \cmd{80} = \rule{2.2cm}{0.3pt};
      \cmd{443} = \rule{2.2cm}{0.3pt}; \cmd{53} = \rule{2.2cm}{0.3pt}.
\hint{22, 80, 443, 53 — match them to SSH, HTTP, HTTPS, DNS.}
\item State the difference between a \textbf{stateless} and a \textbf{stateful} firewall, and say
      which UFW and NSGs are.
\hint{One judges each packet alone; one remembers active connections.}
\item In a firewall rule, name the five things a rule typically specifies.
\hint{Action, protocol, source, destination, port.}
\item What is the security benefit of \textbf{NAT}, even though it is not a firewall?
\hint{It hides internal machines behind one public address.}
\item Explain \textbf{network segmentation} and how it limits ``lateral movement''.
\hint{Separate zones, so a breach in one cannot roam into the others.}
\item True or false: a host firewall is unnecessary if there is a network firewall in front of it.
      Justify.
\hint{Defence in depth — what if the network firewall is bypassed or misconfigured?}
\item What problem does a \textbf{bastion / jump host} solve?
\hint{One hardened, watched entry point instead of exposing every host to SSH.}
\item In Azure, NSG rules are evaluated by \rule{3cm}{0.3pt}; the \rule{2cm}{0.3pt} match wins.
\hint{By priority, lowest number first; the first match wins.}
\item Why does \textbf{SDN} make a cloud network both easier to secure \emph{and} riskier? (one line)
\hint{Central control can program the whole network at once — powerful, and a target.}
\item (Reflection) Which single listening port on your VM would you remove first, and why?
\end{enumerate}

\signoff

\clearpage
\solheader
{\LARGE\bfseries\color{kred} Solutions}\\[2pt]
{\small Work through the lab first, then check here. Model answers are indicative — equivalent reasoning is fine.}\par\vspace{8pt}
\noindent\textbf{Note:} Model answers are indicative -- accept equivalent reasoning. Multiple-choice
answers are in \textbf{bold}.

\wbsection{Knowledge check}
\begin{enumerate}
\item \textbf{An IP address + a port.}
\item Because you cannot anticipate every threat. With default-deny you open only what you can
      justify, so anything unforeseen is blocked automatically -- the principle of least exposure.
\item \cmd{22} = \textbf{SSH}; \cmd{80} = \textbf{HTTP}; \cmd{443} = \textbf{HTTPS}; \cmd{53} =
      \textbf{DNS}.
\item \textbf{Stateless} judges each packet in isolation; \textbf{stateful} tracks active connections
      and automatically allows replies to traffic you started. UFW and NSGs are \textbf{stateful}.
\item \textbf{Action} (allow/deny), \textbf{protocol} (TCP/UDP), \textbf{source}, \textbf{destination},
      and \textbf{port}.
\item NAT hides internal machines behind one public address; they are not directly reachable from
      outside unless a port is explicitly forwarded -- a crude but real first barrier. (Comer Ch.~7,
      p.~91, verified.)
\item \textbf{Segmentation} divides the network into separate zones so a breach in one cannot freely
      reach the others, limiting \emph{lateral movement}. (CSA Domain~7; Comer VLAN/VXLAN p.~89--90.)
\item \textbf{False} -- defence in depth. The host firewall still protects the host if the network
      firewall is bypassed or misconfigured, or if the threat originates \emph{inside} the network.
\item A \textbf{bastion / jump host} provides a single hardened, heavily-monitored entry point for
      admin access (SSH), instead of exposing SSH on every server.
\item Rules are evaluated by \textbf{priority} (lowest number first); the \textbf{first} match wins.
\item SDN centralises and programs the network, so a policy can be applied everywhere at once (easier
      to secure); but the controller is a single high-value target and a misconfiguration propagates
      just as fast (riskier). (Comer Ch.~7, p.~94--95, verified.)
\item Reflection -- no fixed answer; expect students to name an unneeded service bound to a public
      interface.
\end{enumerate}

\wbsection{Lab checkpoints}
\begin{itemize}
\item \cmd{ufw status verbose} should show \emph{default deny (incoming)} plus explicit allow rules for
      22 and 80.
\item Tested from the host: an allowed port responds to \cmd{curl}; a port the firewall drops
      \emph{times out}; with UFW disabled the same closed port answers \emph{connection refused} at once
      -- a good discussion point. Testing against \cmd{localhost} inside the VM proves nothing, because
      UFW does not filter loopback traffic.
\item The before/after \cmd{ss -tulpn} should be unchanged (services still listen); the firewall
      controls \emph{reachability}, not whether a service runs.
\end{itemize}

\wbsection{References}
CSA Security Guidance v5, \textbf{Domain~7: Infrastructure \& Networking} (verified); Comer Ch.~7
(NAT \textbf{p.~91}, VLANs \textbf{p.~89}, VXLAN \textbf{p.~90}, virtual switch \textbf{p.~91}, SDN
\textbf{p.~94}, OpenFlow \textbf{p.~95}) -- all verified; J\o sang Ch.~6 (syllabus, page numbers not
verified). MS Learn: \emph{Describe Azure networking services} (verified, June 2026).

\end{document}
